\section{Discussion}
\label{sec:discussion}

\subsection{What the formal result establishes}

The formal comparison isolates one mechanism. Candidate and parent share the
same five actor tensors, checkpoints, state, and inherited action transform;
only the 60-s completion reserve differs. On ten fresh environment blocks, the
reserve reduced generalized passenger time and improved completion in every
block. The direction also held at every demand level and training seed. This is
strong evidence that the fixed reserve improves this simulator-to-simulator
endpoint relative to its direct parent under the registered scenario family.

The demand pattern is operationally coherent. The benefit is smallest at
demand 0.75 and largest at demand 1.25, where leaving a tail without a follower
has greater passenger consequences. The reserve can add local in-vehicle time,
but it may also prevent a low-load tail bus from advancing in a way that leaves
later demand unserved before the horizon. The formal decrease in unfinished
passengers is consistent with the latter mechanism in these scenarios.
This interpretation is mechanistic rather than causal at the level of
individual passengers because the experiment did not randomize guard
activations within a trajectory.

The result does not establish an active density-ratio benefit. Classifier
validation was near chance, and final training used unit simulator weights.
The learned component should therefore be described as an IQL-initialized,
simulator-continued H2O+ actor with mixed replay, not as a successfully
ratio-corrected policy. The confirmed contribution is the composition of this
fixed actor with the explicit completion-safety reserve.

\subsection{Why development evidence alone was insufficient}

The development history was adaptive. Successive rounds diagnosed actor drift,
hold scale, headway reserve, tail service, and completion, and several rounds
reused the same nine development cells. A favorable result at one stage could
therefore reflect both a real mechanism and accumulated adaptation to those
cells. Earlier variants also used different service margins in some analyses;
passing a looser completion boundary does not imply passing the final 0.003
margin. For these reasons, development improvements were not treated as a
formal effect.

The ten-block confirmation was designed to break that feedback loop. Candidate,
parent, environment cells, endpoints, margin, and tests were frozen before
outcome access. Formal cells were disjoint from training and all development
rounds, and candidate replacement was forbidden. This design, rather than the
number of development iterations, is what supports the bounded confirmed
claim.

\subsection{Limitations}
\label{sec:discussion:limitations}

\paragraph{Data and state representation.}

The offline dataset is large in rows but modest in independent exogenous
variation. Simultaneous decisions and long trajectories create millions of
correlated SMDP tokens from only 30 environment cells. The five training seeds
represent optimizer variation conditional on that fixed data, not five new
datasets. Future data collection should prioritize more independent days,
demand realizations, and disruption patterns instead of denser serialization of
the same cells. The centralized observation also compresses passenger
histories and future service availability. Its Markov sufficiency is a
modeling assumption, not a proved property of the 229-feature encoding.
\Cref{prop:ledger,prop:hurdle} establish exact accounting and action
normalization, not convergence or a completion guarantee.

Positive-action support is also sparse. Exact zero hold dominates replay, and
long holds are uncommon in the final guard states. The completion reserve is
therefore intentionally a transparent physical rule rather than a claim that
the offline actor learned reliable values for 60-s actions in these states. A
future learned completion controller would need explicit tail-state/action
coverage and a constrained or terminal-shortfall objective. Training also uses
total accrued cost, whereas evaluation normalizes each episode by departed
passengers. Their within-cell ordering agrees by \Cref{prop:ledger}, but
cross-demand weighting can differ.

\paragraph{Simulator and runtime scope.}
Training data were collected with SUMO 1.27.1, whereas development and formal
evaluation were locked to SUMO/libsumo 1.25.0. Because formal performance is
measured entirely in 1.25.0, the reported comparison itself is internally
paired. The shift nevertheless complicates attribution between algorithm,
simulator mismatch, and version-specific behavior. A version-sensitivity matrix
with the same fixed policies in both runtimes is needed before claiming
runtime-invariant transfer.

The reduced-order online simulator passed zero-policy and intervention-response
audits, but those tests do not cover every state distribution induced by a
learned policy. It also approximates transfer-service choice through calibrated
frequencies rather than a dynamic first-arrival model. Both environments share
the same timetable, demand source, capacities, and passenger contract, so the
study does not test changes in traveler preferences, fares, dispatch rules, or
fleet composition.

\paragraph{Statistical and comparison scope.}

The exact sign test does not assume normally distributed effect sizes; it
uses independent blocks and a null directional win probability no greater
than one half. It reports directional
consistency, not a confidence interval for a population-average effect. Ten
blocks also provide limited resolution: the smallest attainable one-sided
value is $2^{-10}$. More independent blocks would support effect-size
uncertainty and interaction analyses without relying on episode-level
pseudoreplication.

The lineage-control table is useful context but not a set of confirmatory
pairwise claims. Only the direct parent removes the final reserve while holding
the underlying actor and prior transform fixed. The study is also endpoint-only;
it cannot support a statement about learning speed, AULC, or compute efficiency.
Likewise, the formal result does not retroactively validate unsuccessful
development candidates or the ratio diagnostic.

\paragraph{Operational scope.}

The reserve is interpretable and bounded, but it is not a complete deployment
safety system. The 60-s cap does not encode terminal layovers, driver rules,
accessibility obligations, connection protection, dispatch authority, or
incident response. Before a field trial, the fixed policy should be calibrated
and evaluated on temporally held-out Automatic Vehicle Location, Automatic
Passenger Counting, and fare-collection records, followed by shadow-mode
testing with dispatcher review.

The current evidence ends at a provenance-locked simulator-to-simulator
confirmation. It supports writing and reporting the bounded result, but not a
claim of live passenger benefit. A field study would require separate
authorization, operational stopping rules, and independent evaluation days.
